\begin{enumerate}
\item \hfill(0.5 point for one star, totally 2 points)
\item The field of view of the camera is marked on the map. \hfill(2 p)

% FIGURE NEEDED: solution chart -- the sky map with the CCD image overlaid
% as a tilted square marking the camera's field of view
% (2010_da_sol.pdf page 1, Proceedings of 4th IOAA p. 61)

\item According to the pic of A2, it's easy to find the field of view of the
telescope. It's about $26'$, and the declination of the center of the CCD
image is $7.3^\circ$. Thus the side length of the field of view is:
\[ 26' \times \cos 7.3^\circ = 1550'' \]
Image scale is $d = f\theta$, so, $s = f/206.265$\,mm/arcsec
$= 0.0154$\,mm/arcsec.\\
The chip size is $1550 \times 0.0154 = 24$\,mm. \hfill(4p)

\item The star is 10 pixels across, so the FWHM is $10/3.5 = 2.9$ pixels.
\hfill(4p)\\
Seeing is $S = 2.9\,\mathrm{pixels} \times 1.5''/\mathrm{pixel}$ (from Q3
and 1024 pixels) $= 4.4''$.

\item Theoretical (Airy) diffraction disc is $2.44\lambda/D$ radians in
diameter:
\[ A = 2.44 \times 550 \times 10^{-9}/0.61\ \mathrm{rad} = 0.45'' \sim 0.3\ \text{pixels} \]
$A \ll S$ (seeing). \emph{(Accept all reasonable wavelengths:
450--650\,nm)} \hfill(4p)

\item Seeing $=$ FWHM $\times\ 1.5''/\mathrm{pixel}$ (from Q3) $= 1''$. So,
FWHM $= 1/1.5$ pixel $= 1$ pixel.\\
Printed image of star would then be
$s_2 = 3.5 \times \mathrm{FWHM} = 3.5$ pixels. \hfill(3p)\\
Note: if use: $s_2 = 1'' \times 10\,\mathrm{pix}/4.4'' = 2.3$\,pix, 2 points.

\item For object 1, the trail of the object is about $107''$ (measured from
pic 1, 300\,s exposure). Its angular velocity is:
\[ \omega_1 = 107''/300\,\mathrm{s} = 0.36\,''/\mathrm{s} \]
Note: accept to $v \pm 10\%$. \hfill(3p)

For object 2, it moves about 8 pixels between pic 2A and 2B. 8 pixels
$\sim 12''$, and the time between exposures is $17^{\mathrm{m}}27^{\mathrm{s}}$.
Its angular velocity is: \hfill(3p)
\[ \omega_2 = 12''/1047\,\mathrm{s} = 0.012\,''/\mathrm{s}\ (\text{accept} \pm 10\%). \]
\begin{description}
\item[a)] \textbf{wrong}: different masses of the objects, \hfill($+2$/$-1$p)
\item[b)] \textbf{right}: different distances of the objects from Earth,
  \hfill($+3$/$-1$p)
\item[c)] \textbf{right}: different orbital velocities of the objects,
  \hfill($+3$/$-1$p)
\item[d)] \textbf{wrong}: different projections of the objects' velocities,
  \hfill($+2$/$-1$p)
\item[e)] \textbf{rejected}: Object 1 orbits the Earth while Object 2 orbits
  the Sun. \hfill(0p)
\end{description}
\end{enumerate}
